% This Source Code Form is subject to the terms of the Mozilla Public
% License, v. 2.0. If a copy of the MPL was not distributed with this
% file, You can obtain one at https://mozilla.org/MPL/2.0/.
\documentclass[11pt,a4paper]{article}

\usepackage[margin=2.2cm]{geometry}
\usepackage{amsmath,amssymb,amsthm}
\usepackage{booktabs}
\usepackage{array}
\usepackage{enumitem}
\usepackage{needspace}
\usepackage{microtype}
\usepackage[hidelinks]{hyperref}

% The same source gives the same PDF: no dates, no random identifier and no
% description of the machine that built it. `just pdfs` passes the SHA-256
% of this file as \sourcehash, and the PDF records it, so that
% `just pdfs-check` can tell whether the committed PDF was built from the
% current source.
\ifdefined\pdfinfoomitdate
  \pdfinfoomitdate=1
  \pdftrailerid{}
  \pdfsuppressptexinfo=-1
\fi
\providecommand{\sourcehash}{unknown}
\hypersetup{
  pdftitle={A baseline solution: fixed playout delay},
  pdfsubject={Built from docs/fixed-playout-delay.tex with SHA-256 \sourcehash}
}

\setlist{itemsep=2pt,topsep=4pt,parsep=0pt}
\emergencystretch=3em
\clubpenalty=10000
\widowpenalty=10000

\newcommand{\lean}[1]{\texttt{#1}}
\newcommand{\R}{\mathbb{R}}
\newcommand{\clamp}{\operatorname{clamp}}

\theoremstyle{definition}
\newtheorem{definition}{Definition}[section]
\newtheorem{theorem}[definition]{Theorem}
\newtheorem{remark}[definition]{Remark}

\begin{document}

\begin{center}
{\LARGE\bfseries A baseline solution: fixed playout delay\par}
\end{center}
\medskip

This document describes the first solution to the problem of
\path{docs/specification.pdf}, the baseline that later solutions are
compared with, and what is proved about it. Every speaker plays each frame
of the source a fixed latency $L$ after the frame was captured, as well as it
can tell from its host's clock. The solution is defined in
\path{lean/Unisono/Solutions/FixedDelay.lean}; the notation of the
specification is used throughout.

\section{The solution}

\paragraph{Parameters} (\lean{Params}): the sender's host and the hosts of
the speakers; the nominal rate $\rho$; the latency $L$; the number of copies
$c$ of each message; the period $P$ of the questions for the sender's
clock, in seconds of the asking host's clock; and the slew $w$, the largest
correction of the position per output sample, in frames.

\paragraph{The stream.} When the capture device hands over block $b$, the
sender sends its frames to the host of every speaker other than its own,
with the block's first frame $f(b)$ and its timestamp $\hat c(b)$, a reading
of the sender's clock. Each message is sent $c$ times, so that a link that
drops at most $c - 1$ consecutive messages still delivers one. The copies
leave at the same moment, so they protect against isolated drops only: an
outage, in which a link loses everything sent during some time, takes all
of them. A host keeps, of the blocks that have arrived, the latest: its
first frame $f$ and timestamp $\hat c$.

\paragraph{The clocks.} At a request of one of its devices, a speaker's host
other than the sender's asks the sender for its clock, unless it asked
within the last $P$ seconds of its own clock: it sends its reading $t_1$.
When a question arrives, the sender answers with $t_1$ and its own reading
$t_2$, unless it has already answered a question of that host whose reading
is as late or later; it thus answers each question once, whichever copy
arrives first. When the answer arrives at the host's reading $t_4$, the host
estimates the sender's clock minus its own as
\[
  \hat\theta = t_2 - \frac{t_1 + t_4}{2},
\]
which is NTP's estimate with the sender's two readings taken at the same
moment. The host keeps the estimate of the latest question answered, the
one with the largest $t_1$; answers that arrive out of order do not replace
a newer estimate. Questions and answers are also sent $c$ times. The
sender's host uses $\hat\theta = 0$.

\paragraph{The target} (\lean{target}). When speaker $r$'s device requests a
block whose first output sample is $g$, with timestamp $\hat o$, the host
gives output sample $k$ of the block the target position
\[
  x^*_k = f + \Bigl(\hat o + \frac{k - g}{\rho} + \hat\theta - L - \hat c\Bigr)\,\rho .
\]
The bracket estimates, on the sender's clock, when the sound that output
sample $k$ should play was captured: the time at which the sample leaves,
on the host's clock, converted to the sender's clock, less $L$. The latest
block converts that capture time to a position, at the nominal rate.

\paragraph{The pace} (\lean{advance}, \lean{positions}). A speaker does not
jump to its target. If it chose position $x_{k-1}$ for the output sample
before, it chooses
\[
  x_k = x_{k-1} + 1 + \clamp_w\bigl(x^*_k - x_{k-1} - 1\bigr),
  \qquad \clamp_w(y) = \max\bigl(-w, \min(w, y)\bigr):
\]
one frame on, corrected towards the target by at most $w$ frames. The
position thus advances by between $1 - w$ and $1 + w$ frames per output
sample. The correction is not smoothed: when the target jumps, at a new
block or a new clock estimate, the advance changes by up to $w$ from one
output sample to the next, and by up to $2w$ when the correction reverses.
A speaker continues from the position it chose for its previous output
sample, whether or not the frames for that sample had arrived, so that it
stays on its course through a gap; it starts at its target when it chose
no position there, at its first output sample or while it has no
estimates.

Nothing is retransmitted on request; losses are covered by the copies.

\section{What is proved}

\subsection{Playback at a latency}

The guarantees of a solution follow from one property of its playback. The
file \path{lean/Unisono/Spec/AtLatency.lean} states it for any solution.

\begin{definition}[Playback at a latency]
A playback is \emph{at latency $L$ within $E$ from output sample $K$}
(\lean{World.PlaysAtLatency}) when, for every speaker $r$ and every output
sample $k \ge K$, playback is continuous from $k$ to $k + 1$ and the sound
that $k$ plays was captured within $E$ of $o_r(k) - L$:
$|c(x_{r,k}) - (o_r(k) - L)| \le E$.
\end{definition}

\begin{theorem}\label{thm:atlatency}
Let a playback be at latency $L$ within $E$ from some output sample on, let
$\delta$ bound the time between two output samples from there on, and let
frames be captured at most $\delta_c$ apart. Then, in the long run:
\begin{itemize}
\item it misses no samples and has no discontinuities
  (\lean{missed\_eq\_zero}, \lean{discontinuities\_eq\_zero});
\item its latency is at most $L + E + \delta$ (\lean{meanLatency\_le});
\item its synchronization cost is at most $(E + \delta)^2$
  (\lean{sync\_le}), which is attained when half of the speakers play every
  frame $E + \delta$ early and the other half $E + \delta$ late;
\item its tracking cost is at most $(D + E + \delta)^2$ when every target
  latency is within $D$ of $L$ (\lean{tracking\_le}).
\end{itemize}
\end{theorem}

The proofs bound the cost of each frame and of each output sample from some
point on, and use that a cost accrued at most linearly in time, $B\,T + C$
with $C$ finite, has a long-run average at most $B$ (\lean{longRun\_le}).
The bound on synchronization uses that for values $d_1, \dots, d_R$ within
$M$ of zero, $\sum_{r < r'} (d_r - d_{r'})^2 = R \sum_r d_r^2 - (\sum_r
d_r)^2 \le R^2 M^2$ (\lean{sum\_pairs\_sq\_sub\_le}), Popoviciu's inequality
for the variance.

Playback at a latency says nothing about the sixth objective, roughness:
within $E$ a speaker can change its speed as often as it likes, so a bound on
roughness needs a bound on how a solution changes its speed. A change of
speed by $\Delta v$ between two consecutive steps costs about $\rho\,\Delta
v^2$, since the two steps are about $1/\rho$ apart. The baseline changes its
advance by up to $2w$ at once whenever its target jumps, and its target can
jump at every new block and every new clock estimate, so its roughness is
large; it is not bounded here.

\subsection{The baseline in a bounded world}

\begin{theorem}[The stream]\label{thm:stream}
In a world within the bounds of the bounded model, let the baseline send
each message more times than a link can drop in a row, $c > m_L$. Then every
speaker's host $j$ other than the sender's holds the frames of block $b$ at
every step of the run whose event comes later than
$h(b) + d_{\max}(\text{sender}, j)$ (\lean{block\_arrives}).
\end{theorem}

The theorem relies on the bounded model counting losses per message: the
$c$ copies leave at the same moment and are still consecutive messages on
the link, of which one arrives. A link with outages in time would lose all
of them.

The proof rests on facts about the run of any solution
(\path{lean/Unisono/Spec/Run.lean}): a host keeps the frames it holds; a
message that arrives before the event of a later step has been handled by
then; and the copies that the sender sends to a host at a handoff are its
consecutive messages on that link, of which the bounded model delivers at
least one within $d_{\max}$.

What remains is to show that the baseline's run is at a latency, in every
world within the bounds of the bounded model, with an explicit $E$ and for
every $L$ above an explicit minimum. The proof follows the solution's parts,
each bounding one source of error:
\begin{enumerate}
\item \emph{The views} (proved): a host's view is the list of the events it
  handled, with its clock's readings, and every message of the baseline in
  transit was sent where it should have been (\lean{origin\_inFlight}), so
  what the baseline remembers can be traced to the world. A host other than
  the sender's remembers the latest block that arrived and the estimate of
  the latest question answered (\path{lean/Unisono/Solutions/FixedDelayMemory.lean}).
\item \emph{The clocks} (proved): every question is answered, so at every
  step later than a question plus the longest delays to the sender and back,
  the asking host remembers an estimate from a question as late or later
  (\lean{answer\_reaches}); the estimate it remembers is an NTP sample whose
  question and answer took between the least and the longest delays
  (\lean{offset\_sample}). Such an estimate errs from the true difference of
  the two clocks at the sender's reading by at most half the largest
  difference of the two delays, plus the clocks' rate error times the mean
  delay (\lean{offset\_error}), and the difference of two clocks drifts by
  at most twice their rate error per second (\lean{offset\_drift}).
\item \emph{The target:} the capture time of $x^*_k$ errs from
  $o_r(k) - L$ by at most the clock error, the timestamps' errors
  $2\varepsilon_R$, and the rate errors over the extrapolated spans, so
  $x^*_k$ lies within some $B$ frames of the trajectory $T_k$ whose capture
  time is exactly $o_r(k) - L$.
\item \emph{The pace} (proved): $T_k$ advances by one frame per output
  sample within the devices' rate error $\nu$. When $w > \nu$, a position
  that starts within $B$ of $T_k$, as the baseline's does by starting at
  its target, stays within $B$ of it (\lean{advance\_tracks}).
\item \emph{The data:} the frames that $x_k$ needs, within the resampler's
  lookahead $a$, have arrived when the device requests the block once $L$
  exceeds $E$ plus the longest chain of delays, $d_{\max}$, the handoff
  delay $\eta$, the request lead $\lambda_{\max}$, and the spans of two
  blocks and of the lookahead.
\item \emph{Progress:} the baseline sets no timers, and it sends messages
  only at handoffs and requests, of which finitely many happen before any
  time, and when a question arrives; it answers each question at most once,
  and blocks and answers lead to no message. So finitely many events happen
  before any time, and the run advances past every true time.
\end{enumerate}
Together with Theorem~\ref{thm:atlatency}, this gives the baseline's
guarantees under the bounded model.

The estimate $\hat\theta$ errs by up to half the largest difference of the
delays of the question and the answer. When every link's delays lie within
the same $[d_{\min}, d_{\max}]$, that is $(d_{\max} - d_{\min})/2$, so two
speakers can be $d_{\max} - d_{\min}$ apart. When the two directions of a
link have different bounds, the midpoint carries a bias of up to half the
difference of the two delays even when they do not vary, and no exchange of
this kind can detect it. With the same bounds on every link, the theorem of
Lundelius and Lynch on clock synchronization says that no algorithm keeps
the clocks of $n$ hosts closer than $(1 - 1/n)(d_{\max} - d_{\min})$. Hosts
that all exchange readings with each other attain it, which gains most when
there are few hosts: with three, the limit is two thirds of
$d_{\max} - d_{\min}$, while exchanges with the sender alone, as in the
baseline, allow all of it. Carried over to playback, that is the limit the
baseline will be compared with.

A pause, in which a link holds messages back and releases them later,
enters the bounded model only through $d_{\max}$. A link that holds messages
for tens of milliseconds therefore makes the baseline's clock estimate err
by tens of milliseconds. Keeping the estimates of the exchanges with the
shortest round trip, as NTP's clock filter does, avoids most of that error,
and is the first improvement on the baseline that this suggests.

\end{document}
